% Page 057

\seite{57}

\abschnitt{Schuljahr 1921–22}
\pstart
\margmark{Entlassen wu\\die Landwirt\\Schüler auf}%
rden in diesem Jahre 8 Kinder, welche sich für\ob
schaft entschlossen, nämlich Schüler: 1\ob
Ostern 1 Schüler und Knaben und 3 Mädchen; 15 Mädchen

\pend
\abschnitt{Lehrerwechsel}
\pstart

\margmark{Am 1. Ap\\hiesiger Eig\\Karolich [?]\\Schule wurde}%
ril wurde der Lehrerin Platz in\ob
enschaft in die umfangreichere Schule nach\ob
 (bei Daun) versetzt.  Für die hiesige Stelle\ob
 der Schulamtsbewerber Veit ernannt.

\margmark{Nachdem die\\wärtiger Tät\\15. Mai 1921  S\\bewerber Vei\\des 1. Mai d\\7. April 192\\der Schule i\\aufgestellte\\Möge zum Woh\\meinde gedei}%
Reform- Diät und Gehaltsangelegenheiten gegen\ohb
igkeit die hiesige Schaffung, infolge der neuen\ob
timmungen verlassen mußte, wurde der Schulamts\ohb
t, nach der Verordnung vom 1.\ April\ob
ienstfrei gelaßt. Laut Regierungsbeschluß vom\ob
1 wurde Lehrer Kneupp, seit 28.\ 5. 19\unsicher{} an\ob
n Flußen\unsicher, auf den 1.\ Mai als nunmehriger\ob
r Lehrer an die hiesige Stelle berufen.\ob
lsein\unsicher{} für Schule und Ge\ohb
hlich sein!

\margmark{Am 13. M\\Unsre Gemein\\11 Mädchen,\\Schule schon\\und feierlic\\in festem Ga\\Reden nebstd\\Firsten und}%
ai fand in festlicher Stimmung statt.\ob
de umfaßte 15 Jünglinge,\ob
4 Renten\unsicher, von denen etwa 2 aus der\ob
 entlassen waren. Festlich geschmückt\ob
h begann\unsicher{} geradewegs nach seiner Weise\unsicher\ob
nge zum Orte, über die Geistlichen\ob
em bei der Andacht die betreffenden\ob
Geistlichen\unsicher.

\margmark{24. Mai  Vom 14\\die achttägi\\schönem Rege}%
. bis einschließlich 23.\ Mai dauerten\ob
gen Pfingstferien, die von günstig\ob
n begleitet waren.
\pend
